% Figure: completing the square rendered as a geometric area-balance.
% A square of side x plus a strip of width b/2 on two sides equals a
% larger square minus the corner square of area (b/2)^2.
\begin{figure}[htbp]
\centering
\begin{tikzpicture}[
  scale=1.0,
  fillA/.style={fill=engblue!18, draw=engblack, thick},
  fillB/.style={fill=engblue!8, draw=engblack, thick},
  fillC/.style={fill=engred!18, draw=engblack, thick},
  ann/.style={font=\small, align=center},
]
% Left panel: x^2 + b x laid out as a square plus two strips.
\begin{scope}[shift={(0,0)}]
  % central x^2 square
  \draw[fillA] (0,0) rectangle (3,3);
  \node[ann] at (1.5,1.5) {$x^{2}$};
  % strip to the right (width b/2)
  \draw[fillB] (3,0) rectangle (4,3);
  \node[ann] at (3.5,1.5) {$\tfrac{b}{2}x$};
  % strip on top (width b/2)
  \draw[fillB] (0,3) rectangle (3,4);
  \node[ann] at (1.5,3.5) {$\tfrac{b}{2}x$};
  % side labels
  \node[ann] at (1.5,-0.4) {$x$};
  \node[ann] at (3.5,-0.4) {$\tfrac{b}{2}$};
  \node[ann] at (-0.45,1.5) {$x$};
  \node[ann] at (-0.45,3.5) {$\tfrac{b}{2}$};
  \node[ann] at (2.0,-1.2) {$x^{2}+bx$};
\end{scope}
% Right panel: completed square (x+b/2)^2 minus the missing corner.
\begin{scope}[shift={(6.5,0)}]
  \draw[fillA] (0,0) rectangle (3,3);
  \node[ann] at (1.5,1.5) {$x^{2}$};
  \draw[fillB] (3,0) rectangle (4,3);
  \node[ann] at (3.5,1.5) {$\tfrac{b}{2}x$};
  \draw[fillB] (0,3) rectangle (3,4);
  \node[ann] at (1.5,3.5) {$\tfrac{b}{2}x$};
  % the corner that completes the big square
  \draw[fillC] (3,3) rectangle (4,4);
  \node[ann] at (3.5,3.5) {$\bigl(\tfrac{b}{2}\bigr)^{2}$};
  \node[ann] at (2.0,-0.5) {$\bigl(x+\tfrac{b}{2}\bigr)^{2}$};
  \node[ann] at (2.0,-1.2)
    {$x^{2}+bx=\bigl(x+\tfrac{b}{2}\bigr)^{2}-\bigl(\tfrac{b}{2}\bigr)^{2}$};
\end{scope}
\end{tikzpicture}
\caption[Completing the square as an area balance]{Completing the
square is an area balance. The expression $x^{2}+bx$ tiles a square of
side $x$ together with two strips of width $b/2$. Adding the red corner
of area $(b/2)^{2}$ completes a square of side $x+b/2$, so $x^{2}+bx =
(x+b/2)^{2}-(b/2)^{2}$. The identity is the algebraic engine behind the
quadratic formula, the minimum of a parabola, and the Gaussian
integral the reader meets in Chapter 6.}
\label{fig:vol02:ch01:completing-square}
\end{figure}
